\hnsection{FREE LIE ALGEBRAS}

It remains to speak of a series of works on Lie algebras where the connection with the theory of Lie groups is very tenuous; this research has on the other hand important applications to the theory of ``abstract'' groups and more especially nilpotent groups.

Its origin is the work of P. Hall {[}24{]}, which appeared in 1932. However Lie algebras are not discussed here: P. Hall had in mind the study of a certain class of p-groups, those which he called ``regular''. But this led him to examine in detail iterated commutators and the lower central series of a group; on this subject he established a version of the Jacobi identity (cf.~Chapter II, § 4, no. 4, formula (20)) and the ``Hall formula''

\[
(x y) ^ {n} = x ^ {n} y ^ {n} (x, y) ^ {n (1 - n) / 2} \dots \quad (\text { cf.   Chapter   II }, \S 5, \text { Exercise   9 }).
\]

Almost immediately (in 1935--1937) appeared the fundamental works of W. Magnus ({[}25 a{]} and {[}25 b{]}) and E. Witt {[}30{]}. In {[}25 a{]} Magnus used the same algebra of formal power series \(\hat{\mathbf{A}}\) as Hausdorff (since called ``Magnus algebra''); he embedded the free group F in it and used the natural filtration of \(\hat{\mathbf{A}}\) to obtain a decreasing sequence \((\mathbf{F}_{n})\) of subgroups of F; it is one of the first examples of a filtration. He conjectured that the \(F_{n}\) coincide with the terms of the lower central series of F. This conjecture was proved in his second memoir {[}25 b{]}; also in this work he showed explicitly the close relation between his ideas and those of P. Hall and defined the free Lie algebra L (as a subalgebra of \(\hat{\mathbf{A}}\)) which he showed essentially to be identified with the graded algebra of F. In {[}30{]}, Witt completed this result on various points. He showed notably that the enveloping algebra of L is a free associative algebra and deduced immediately the rank of the homogeneous components of L (``Witt formulae'').

As for the proof of the basis of L known as the ``Hall basis'' (cf.~Chapter II, § 2, no. 11), it seems that it only appeared in 1950 in a note of M. Hall {[}40{]}, although it was implicit in the works of P. Hall and W. Magnus quoted above.

